\documentclass[11pt,a4paper]{article}
\usepackage[margin=22mm]{geometry}
\usepackage[T1]{fontenc}
\usepackage{lmodern}
\setlength{\parindent}{0pt}
\setlength{\parskip}{6pt}
\makeatletter
\renewcommand\section{\@startsection{section}{1}{0pt}{10pt}{4pt}{\normalfont\normalsize\bfseries}}
\makeatother
\pagestyle{empty}
\begin{document}
\begin{center}
{\LARGE\bfseries CampusPool\par}
\vspace{3pt}
{\large Campus Ride and Activity Coordination Website\par}
\vspace{7pt}
UCS503P Software Engineering Project Proposal\\
Thapar Institute of Engineering and Technology, Patiala\\
Academic year 2026--27
\vspace{10pt}

\begin{tabular}{@{}lr@{}}
\textbf{Team member} & \textbf{Roll number}\\
Pihu & 1024031014\\
Jayant Kaushik & 1024031016\\
Aashna & 1024031018
\end{tabular}
\end{center}

\section{Problem and objective}
Students coordinate shared rides and campus activities through group chats, where posts, seat availability and confirmations can be difficult to track. Participants must also calculate and follow up on shared ride expenses. CampusPool proposes a campus-focused website to bring plan discovery, participation management and fare tracking into one place.

\section{Proposed solution}
Users will sign in through a one-time code sent to their Thapar email address. They can publish rides or activities, discover suitable plans and request to join. Organisers will accept or decline requests, and participants will be able to view their status.

After a ride, the organiser will enter the actual fare. The website will calculate equal shares across all participants, including the organiser, and provide UPI payment links. Riders will acknowledge payment, and the organiser will confirm receipt. Non-travel activities will follow the participation workflow without fare settlement. Chat, notifications and reporting will support coordination and moderation.

\section{Technical approach}
The website will use Next.js and TypeScript for the frontend, FastAPI for the backend and PostgreSQL for data storage. SQLAlchemy and Alembic will support database access and migrations. WebSockets will deliver live updates. Docker and Render will support deployment. Backend validation and access controls will protect organiser actions and user records.


\section{Prototype and delivery plan}
The first prototype will demonstrate the complete ride-request workflow between two accounts. One user publishes a ride, another discovers it and requests a seat, and the organiser accepts or declines. Both accounts should display the resulting status. This provides an early demonstration of the core coordination problem.

Later increments will add non-travel activities, fare splitting, payment acknowledgement, chat, notifications and moderation. Each increment will be reviewed against its requirements before integration. The GitHub repository will hold application code, design documentation, reports and individual weekly journals.

\newpage
\section{Evaluation criteria}
Evaluation will combine functional checks with a short user walkthrough.
\begin{itemize}
\item \textbf{Workflow completion.} Users should be able to publish a ride, find it, submit a request and view the organiser's decision. Record completion time and any step requiring assistance.
\item \textbf{Capacity and permissions.} Check that accepted participants stay within the seat limit, including competing requests for the final seat. Confirm that only the organiser can decide requests and confirm payments for that ride.
\item \textbf{Fare accuracy.} Compare calculated shares with known examples. For an actual fare of INR~600 shared by four participants, including the organiser, each share should be INR~150.
\item \textbf{Input and account validation.} Check missing fields, invalid email addresses and invalid or expired one-time codes. Repeat the activity workflow and confirm that it creates no fare records.
\end{itemize}

\section{Scalability and available tools}
Next.js, FastAPI and PostgreSQL provide established components for the interface, API and persistence. SQLAlchemy, Alembic and pytest support database access, migrations and backend testing. These tools allow the team to focus on coordination workflows and correctness.

The initial deployment will use a single backend process. As usage grows, pagination and suitable database indexes will support discovery queries. Multiple backend workers will require shared event delivery and rate-limit state so that connected users receive consistent updates.

\section{Risks and mitigation}
Email delivery delays can interrupt sign-in, so OTP delivery will be checked on the hosted build and the interface will provide retry feedback. Server-side validation and database transactions will protect participation records. Environment variables will hold deployment secrets, and synthetic accounts will be used during testing.

Payment records depend on rider acknowledgement and organiser confirmation. The interface will distinguish these steps clearly. Hosting and database configuration will be checked through health checks and complete user walkthroughs before evaluation.

\section{Project deliverables}
The project will deliver the CampusPool website with its frontend, backend and database, a GitHub source repository, UML and data flow diagrams, a deployment configuration, project reports, individual weekly journals and a final demonstration.

\end{document}
